\medskip
\noindent
13. \underline{Links with transcendental theory: the complex analytic
case}.

For the general notions concerning analytic spaces, we refer to [2].

\medskip
\noindent
13.1. \underline{Review of analytic groups over $S$.} Let $S$ be an
analytic space, and let $A$ be an analytic space over $S$, endowed with
the structure of a commutative group over $S$ (that is, $A$ is a
commutative group in the category of analytic spaces over $S$). We assume
that $A$ is \underline{smooth} over $S$, and denote by

\begin{equation}\tag{13.1.1}
 \underline E=\underline{\operatorname{Lie}}(A/S)
\end{equation}

\noindent
the locally free module on $S$ obtained by pulling the relative tangent
sheaf of $A$ over $S$ back along the zero section. We shall admit here the
definition of the exponential homomorphism

\begin{equation}\tag{13.1.2}
 \exp:E\longrightarrow A\qquad ,
\end{equation}

\noindent
where $E$ is the analytic vector bundle on $S$ whose module of
holomorphic sections is $\underline E$. In a neighbourhood of a point
$s\in S$, this homomorphism is first defined near the zero section of $E$
by the familiar methods of the theory of formal groups in characteristic
zero, and is then extended in the evident way, by multiplicativity, to
all of $E$. It is a homomorphism of relative analytic groups over $S$;
it is \underline{\'etale}, and its image is the open subspace $A^o$ of
$A$ which is the union of the identity components of the fibres of $A$
over $S$. If $R$ is the analytic subspace which is the kernel of $\exp$,
we therefore have an exact sequence of relative analytic groups

\begin{equation}\tag{13.1.3}
 0\longrightarrow R\longrightarrow E\longrightarrow A^o
 \longrightarrow0\qquad .
\end{equation}

\noindent
Here the morphism $E\longrightarrow A^o$ is surjective and \'etale, so
$R$ is a relative analytic group \underline{\'etale} over $S$. It is an
analytic subspace of $E$, and is closed in $E$ if and only if $A^o$ is
separated over $S$, that is, if and only if the zero section of $A^o$ is
closed in $A^o$.

\medskip
\noindent
13.1.4. Of course, the exact sequence (13.1.3) is
\underline{functorial in $A$}, and commutes with every base change
$S'\longrightarrow S$. In this way one obtains an equivalence of
categories, compatible with base changes, between the category of
relative analytic groups over $S$ which are smooth and have connected
fibres, and the category of pairs $(E,R)$ consisting of a locally free
vector bundle $E$ over $S$ and an analytic subgroup $R$ of $E$ which is
\'etale over $S$.

\medskip
\noindent
13.1.5. When $A$ is an ``\underline{abeloid}'' group over $S$, by which
we mean that it is smooth, proper, and separated over $S$, with connected
fibres---so that in particular $A=A^o$---then $R$ is a locally constant
group over $S$, whose rank at each point $s$ (as a $\mathbb Z$-module) is
equal to $2\operatorname{rank}_s(E)$, and conversely. In this way one
obtains the familiar interpretation of an abeloid space over $S$ in
terms of a ``\underline{local system}'' $R$ over $S$ and a
\underline{complex} analytic vector-bundle structure on the real
analytic bundle $R\otimes_{\mathbb Z}\mathbb R$ defined by $R$.

\medskip
\noindent
13.2. \underline{Formal completion along a fibre.} We now fix a point
$s\in S$, and shall work locally in a neighbourhood of $s$. For every
integer $n\geqslant0$, let $S_n$ denote the $n$-th infinitesimal
neighbourhood of $s$ in $S$, and, for every analytic space $X$ over $S$,
put likewise $X_n=X\times_S S_n$. As $n$ varies, the system of the $X_n$
then forms an inductive system of analytic spaces over $S$, which we
denote by $\widehat X$ and call the \underline{formal completion of $X$
along $X_s$}. It clearly depends functorially on $X$, and its formation
commutes with finite projective limits; hence $X\longmapsto\widehat X$
takes groups to groups. Notice that $\widehat X$ in fact depends only on
the restriction of $X$ to an arbitrary open neighbourhood $U$ of $s$ in
$S$.

\medskip
\noindent
\underline{Lemma 13.2.1.} \underline{The data are those of 13.1, and
suppose that the following conditions hold:} a) $A=A^o$;
b) $R$ \underline{is constant in a neighbourhood of $s$}; and
c) $R_s$ \underline{spans the complex vector space $E_s$}.
\underline{Let $A'$ be a second smooth analytic group over $S$, and
consider the map}

\begin{equation*}
 u\longmapsto\widehat u:
 \operatorname{Hom}_{S\text{-}\mathrm{gr}}(A,A')
 \longrightarrow
 \operatorname{Hom}_{\widehat S\text{-}\mathrm{gr}}
   (\widehat A,\widehat {A}')\quad ,
\end{equation*}

\noindent
\underline{where the notation is that of 13.2, and the analogous map
obtained by restricting over a variable open neighbourhood $U$ of $s$
and passing to the inductive limit:}

\begin{equation}\tag{13.2.1.1}
 u\longmapsto\widehat u:
 \operatorname{Hom}_{(S,s)\text{-}\mathrm{gr}}(A,A')
 \longrightarrow
 \operatorname{Hom}_{\widehat S\text{-}\mathrm{gr}}
   (\widehat A,\widehat {A}')\quad .
\end{equation}

\noindent
(\underline{Here $(S,s)$ denotes the germ of the analytic space $S$ at
$s$.}) \underline{Then the latter map is bijective.}

Indeed, consider the exact sequence

\begin{equation*}
 0\longrightarrow R'\longrightarrow E'\longrightarrow {A'}^o
 \longrightarrow C
\end{equation*}

\noindent
analogous to (13.1.3) and defined by $A'$. Since the homomorphisms
$A\longrightarrow A'$ necessarily factor through ${A'}^o$, and the same
holds after the base changes $U\longrightarrow S$ and
$S_n\longrightarrow S$, we may assume that $A'={A'}^o$. By the
dictionary of 13.1.4, it remains to prove that giving a germ of a vector
bundle homomorphism $E\longrightarrow E'$ which takes $R$ into $R'$ is
equivalent to giving a coherent system of homomorphisms
$E_n\longrightarrow E'_n$ which take $R_n$ into $R'_n$.

If $\underline{\mathcal O}_s$ denotes the local ring of $S$ at $s$,
giving a coherent system of homomorphisms
$E_n\longrightarrow E'_n$ is equivalent to giving a homomorphism of
$\widehat{\underline{\mathcal O}}_s$-modules

\begin{equation*}
 \widehat F:\widehat{\underline E}_s
 \longrightarrow\widehat{\underline E}'_s,
\end{equation*}

\noindent
where $\underline E_s$ and $\underline E'_s$ are the free
$\underline{\mathcal O}_s$-modules obtained at $s$ from the
$\underline{\mathcal O}_s$-modules $\underline E$ and $\underline E'$,
corresponding to $E$ and $E'$. On the other hand, a germ of a vector
bundle homomorphism $E\longrightarrow E'$ is equivalent to giving a
homomorphism of $\underline{\mathcal O}_s$-modules

\begin{equation*}
 F:\underline E_s\longrightarrow\underline E'_s\quad .
\end{equation*}

\noindent
This already shows that (13.2.1.1) is injective. For surjectivity, it is
enough to prove that if $\widehat F$ is a homomorphism as above such that,
for every $n$, the induced homomorphism
$F_n:E_n\longrightarrow E'_n$ takes $R_n$ into $R'_n$, then
$\widehat F$ comes from an $F$---that is, it takes $\underline E_s$ into
$\underline E'_s$---and the corresponding germ of homomorphisms
$E\longrightarrow E'$ takes $R$ into $R'$.

By hypothesis c), $R_s=R_o$ spans $E_s$; it follows from Nakayama's
lemma that the group $R_o$ (identified either with the group of sections
of $\widehat R$ over $\widehat S$, or with the group of germs of sections
of $R$ over $S$) generates $\widehat E_s$. Thus, to prove that
$\widehat F(\underline E_s)\subset\underline E'_s$, it is enough to prove
that $\widehat F(R_o)\subset\underline E'_s$. But
$\widehat F(R_o)\subset R'_o$, where $R'_o$ is itself interpreted as the
group of germs of sections of $R'\subset E'$ in a neighbourhood of $s$.
This gives the asserted inclusion, and hence the existence of $F$.
Moreover, we see that $F(R_o)\subset R'_o$, where $R_o$ and $R'_o$ are
regarded as groups of germs of sections of $E$ and $E'$, respectively.
By b), $R_o$ generates $R$ in a neighbourhood of $s$; it follows that
$F(R)\subset R'$ in a neighbourhood of $s$, which completes the proof.

\medskip
\noindent
\underline{Corollary 13.2.2.} \underline{If $A$ satisfies conditions
a), b), and c) of 13.2.1, it is determined up to unique isomorphism in a
neighbourhood of $s$ by the knowledge of the $\widehat S$-group
$\widehat A$.}\footnote{The printed source refers here to the conditions
of 13.2.2 itself; the three conditions are introduced in Lemma 13.2.1.}
